\documentclass[letterpaper,11pt]{article}

\usepackage[T1]{fontenc}
\usepackage[utf8]{inputenc}
\usepackage[bitstream-charter]{mathdesign}
\usepackage[top=0.17in,bottom=0.22in,left=0.38in,right=0.38in]{geometry}
\usepackage{titlesec}
\usepackage{enumitem}
\usepackage[hidelinks]{hyperref}
\usepackage{fancyhdr}
\usepackage[english]{babel}
\usepackage{tabularx}
\usepackage{array}
\usepackage[protrusion=true,expansion=false]{microtype}
\ifdefined\pdfgentounicode\input{glyphtounicode}\fi

\pagestyle{fancy}
\fancyhf{}
\renewcommand{\headrulewidth}{0pt}
\renewcommand{\footrulewidth}{0pt}
\setlength{\footskip}{12pt}

\urlstyle{same}
\raggedright
\setlength{\tabcolsep}{0in}
\linespread{0.97}
\ifdefined\pdfgentounicode\pdfgentounicode=1\fi

%----------SECTION STYLE: bold serif title + gap + full-width rule below----------
\titleformat{\section}
  {\large\bfseries}
  {}{0em}{}
  [\vspace{2pt}\titlerule]
\titlespacing*{\section}{0pt}{4pt}{2pt}

%----------LIST STYLE: small bullet, hanging indent, target line height----------
\renewcommand\labelitemi{$\vcenter{\hbox{\scriptsize$\bullet$}}$}
\setlist[itemize]{
  leftmargin=0.18in,
  labelsep=0.08in,
  itemsep=0pt,
  topsep=1.5pt,
  parsep=0pt,
  partopsep=0pt
}

%----------ENTRY COMMANDS----------
% Education: bold degree + school left, dates right
\newcommand{\eduItem}[3]{%
  \begin{tabular*}{\textwidth}{@{}l@{\extracolsep{\fill}}r@{}}
    \textbf{#1}, #2 & #3 \\
  \end{tabular*}\par
}

% Experience: Company (bold), Role, | Location (italic) left; Dates right
\newcommand{\expItem}[4]{%
  \vspace{2pt}
  \begin{tabular*}{\textwidth}{@{}l@{\extracolsep{\fill}}r@{}}
    \textbf{#1}, #2 \textit{| #3} & #4 \\
  \end{tabular*}\par\vspace{-3pt}
}

% Project: name (bold) left, tech stack right, same line
\newcommand{\projItem}[2]{%
  \vspace{2pt}
  \begin{tabular*}{\textwidth}{@{}l@{\extracolsep{\fill}}r@{}}
    \textbf{#1} & #2 \\
  \end{tabular*}\par\vspace{-3pt}
}

% Skills: bold label run-in
\newcommand{\skillLine}[2]{\textbf{#1:} #2\par\vspace{1.5pt}}

% Compact inline block with slight left indent (publications, open source & community)
\newcommand{\inlineBlock}[1]{{\setlength{\leftskip}{0.18in}\small #1\par}}

\begin{document}

%----------HEADING----------
\begin{center}
    {\fontsize{18}{20}\selectfont\bfseries DEVA ANAND}\\ \vspace{3pt}
    {\small
    \href{mailto:devaanand@umass.edu}{devaanand@umass.edu}
    \enspace$\bullet$\enspace
    (413) 315-1715
    \enspace$\bullet$\enspace
    \href{https://devaanand.com}{devaanand.com}
    \enspace$\bullet$\enspace
    \href{https://linkedin.com/in/devaaa/}{linkedin.com/in/devaaa}
    \enspace$\bullet$\enspace
    \href{https://github.com/Deva-1903}{github.com/Deva-1903}}\\ \vspace{2pt}
    {\small\textit{Available for Fall 2026 internship: September 14 -- December 30, 2026 (full-time)}}
\end{center}
\vspace{-8pt}

%----------EDUCATION----------
\section{Education}
\eduItem{MS, Computer Science}{University of Massachusetts Amherst}{2025--2027}
{\small\textit{Advanced NLP, Advanced Machine Learning, Optimization in CS, Systems for Data Science, Performance Engineering}}\par
\vspace{3pt}
\eduItem{BE, Computer Science and Engineering}{Vels Institute of Science, Technology and Advanced Studies}{2019--2023}

%----------EXPERIENCE----------
\section{Professional Experience}

\expItem{Culture \& Morality Lab, UMass Amherst}{Lab Assistant}{Amherst, MA}{Jul 2026 -- Present}
\begin{itemize}
  \item Sole engineer building the lab's CCR (Contextualized Construct Representations) text-analysis platform (Python, \textbf{FastAPI}, sentence-transformers, React), productionizing a published NLP method into a self-serve embedding-and-scoring service for researchers.
  \item Built a YAML model registry, a versioned append-only construct library (\textbf{99 constructs from 38 questionnaires}), and a structured data-quality warning system with deterministic language detection; grew the backend suite to \textbf{40 hermetic tests} that run in CI without ML dependencies.
\end{itemize}

\expItem{\href{https://www.wysa.io}{Wysa}}{Backend Engineer}{Bengaluru, India}{Jul 2023 -- Jul 2025}
\begin{itemize}
  \item Engineered a multi-tenant NHS eTriage backend (Node.js, MongoDB) serving \textbf{8+ UK clinical clients} and \textbf{10{,}000+ monthly triage submissions} through REST integrations with Mayden's iaptus platform; contributed to \textbf{\$340K+} in revenue.
  \item Architected a full-stack training-data annotation platform (React, Node.js, MongoDB) used daily by the AI team, building the pipeline that generated labeled training examples and cutting manual annotation by \textbf{100+ hours/month} while lifting labeling throughput by \textbf{60\%}.
  \item Diagnosed and resolved production issues across tenants, monitoring logs during releases and client integrations and shipping targeted fixes for tenant-specific incidents to keep the multi-tenant eTriage service stable and reliable.
  \item Owned backend security and reliability across the eTriage stack, implementing \textbf{AWS KMS} encryption for clinical PII, hardening authentication and rate limiting, and remediating \textbf{20+ VAPT findings} with regression-tested fixes.
\end{itemize}

\expItem{Cario Growth Services}{Backend Developer Intern}{Chennai, India}{Nov 2022 -- May 2023}
\begin{itemize}
  \item Designed and implemented \textbf{20+ RESTful APIs} in Node.js + Fastify backed by PostgreSQL with input validation, clear contracts, and backward-compatible changes used in production by the frontend team.
  \item Shipped a production LLM auto-commenting service using \textbf{4-bit quantized Llama-7B} via llama.cpp on a single VM, with async queued generation, per-character rate limits, short-circuit caching, and a deterministic content-safety filter.
\end{itemize}

%----------PROJECTS----------
\section{Projects}

\projItem{Spark ETL Performance Optimization {\normalfont\small(\href{https://github.com/Deva-1903/Spark-ETL-Optimization}{GitHub})}}{PySpark, Spark UI, SQL, NYC TLC Dataset}
\begin{itemize}
  \item Cut end-to-end Spark ETL runtime \textbf{$\sim$25\%} (17--19 min to 14.3 min) on a \textbf{$\sim$1.4B-row} NYC TLC dataset by swapping shuffle joins for broadcast joins (minutes to $\sim$0.05s), enabling AQE and skew-join handling, and tuning shuffle partitions, schema-aware column-pruned reads, and year/month partitioning.
\end{itemize}

\projItem{Cache- and SIMD-Aware Matrix Multiplication {\normalfont\small(\href{https://github.com/Deva-1903/cs690pf}{GitHub})}}{C++, AVX/SIMD, OpenMP, Linux perf}
\begin{itemize}
  \item Optimized C++ matrix multiplication from a naive baseline through cache-aware tiling, manual \textbf{AVX/SIMD} vectorization, register-aware micro-kernels, and \textbf{OpenMP} multithreading, profiling L1/L3/TLB misses with \textbf{perf} and inspecting generated assembly to localize memory-hierarchy bottlenecks.
\end{itemize}

\projItem{KG2RAG-Enhanced -- Multi-Hop QA on HotpotQA {\normalfont\small(\href{https://github.com/Deva-1903/KG2RAG-685-NLP}{GitHub})}}{Python, Ollama / LLaMA-3, sentence-transformers, RRF}
\begin{itemize}
  \item Built the multi-view retrieval and ranking component of a KG-guided RAG pipeline (query-side decomposition, RRF fusion, cross-encoder reranking, MMR), lifting supporting-fact recall \textbf{+2.68\%} (54.53 to 57.21) over the KG\textsuperscript{2}RAG baseline on a 4{,}905-question evaluation.
\end{itemize}

%----------PUBLICATIONS----------
\section{Publications}
\vspace{2pt}
\inlineBlock{- \textit{``Alzheimer's Disease Classification using Transfer Learning,''} IEEE CONIT 2023, first author, deep transfer learning on neuroimaging data. {\normalfont\small(\href{https://ieeexplore.ieee.org/document/10205760}{Paper})}}

%----------SKILLS----------
\section{Skills}
\skillLine{Languages}{Python, C++, Java, Go, TypeScript, JavaScript, SQL, Bash}
\skillLine{Systems \& Data}{Spark / PySpark, Hadoop, Kafka, distributed systems, load balancing, Linux perf, SIMD, OpenMP}
\skillLine{Backend}{Node.js, Express, FastAPI, REST APIs, async queues, MongoDB, PostgreSQL, Redis, DynamoDB}
\skillLine{ML \& Retrieval}{PyTorch, sentence-transformers, embeddings, RRF, cross-encoder reranking, MMR, RAG, llama.cpp}
\skillLine{Cloud \& DevOps}{AWS, GCP, Azure, Kubernetes, Docker, Git, GitHub Actions, CI/CD}

\end{document}
